\RequirePackage{iftex}
\ifPDFTeX
\documentclass[a4paper,12pt]{article}
\usepackage[margin=22mm]{geometry}
\usepackage[T1]{fontenc}
\usepackage[utf8]{inputenc}
\usepackage{graphicx}
\usepackage{CJKutf8}
\else
\documentclass[a4paper,12pt]{jsarticle}
\usepackage[dvipdfmx,margin=22mm]{geometry}
\usepackage[dvipdfmx]{graphicx}
\fi
\usepackage{amsmath,amssymb,amsthm}
\usepackage{enumitem}

\setlength{\parindent}{0pt}
\setlength{\parskip}{0.35em}
\setlength{\fboxsep}{8pt}
\setlist[itemize]{leftmargin=2em,itemsep=0.25em,topsep=0.35em}
\setlist[enumerate]{leftmargin=2.2em,itemsep=0.45em,topsep=0.35em}

\newcommand{\pointbox}[1]{%
\begin{center}
\fbox{%
\begin{minipage}{0.91\linewidth}
#1
\end{minipage}}
\end{center}
}

\begin{document}
\ifPDFTeX
\begin{CJK}{UTF8}{min}
\fi

\theoremstyle{definition}
\newtheorem*{definition}{定義}
\newtheorem*{theorem}{定理}
\newtheorem*{example}{例}
\newtheorem*{remark}{注意}
\renewcommand{\proofname}{\normalfont 証明}

\begin{center}
{\LARGE 数学1A 第10回レジュメ}\\[0.4em]
実数の連続性・数列の極限
\end{center}

\section*{1. 上限と実数の連続性}

\begin{definition}
$A\subset\mathbb{R}$ とする。実数 $M$ が $A$ の上界であるとは、任意の
$a\in A$ について $a\le M$ が成り立つことをいう。

$A$ の上界全体の中で最小のものを $A$ の上限といい、
\[
\sup A
\]
と書く。
\end{definition}

\begin{theorem}[上限の近似性]
$A$ が空でなく上に有界で、$\alpha=\sup A$ とする。このとき任意の
$\varepsilon>0$ に対して、ある $a\in A$ が存在して
\[
\alpha-\varepsilon<a\le \alpha
\]
となる。
\end{theorem}

\begin{proof}
$\alpha$ は上界なので、任意の $a\in A$ について $a\le\alpha$ である。
もし $\alpha-\varepsilon<a$ となる $a\in A$ が存在しなければ、すべての
$a\in A$ について $a\le \alpha-\varepsilon$ となる。したがって
$\alpha-\varepsilon$ は $A$ の上界である。これは $\alpha$ が最小の上界である
ことに矛盾する。
\end{proof}

\pointbox{
上限 $\sup A$ は、集合 $A$ の元そのものとは限らない。しかし、$A$ の元で
上限にいくらでも近いものは必ず存在する。
}

\section*{2. アルキメデスの性質}

\begin{theorem}[アルキメデスの性質]
任意の $x\in\mathbb{R}$ に対して、ある $n\in\mathbb{N}$ が存在して
\[
n>x
\]
となる。
\end{theorem}

\begin{proof}
もし任意の $n\in\mathbb{N}$ について $n\le x$ ならば、$\mathbb{N}$ は上に有界で
ある。そこで $\alpha=\sup\mathbb{N}$ とおく。上限の近似性より、ある
$m\in\mathbb{N}$ が存在して $\alpha-1<m$ となる。すると $m+1>\alpha$ である。
しかし $m+1\in\mathbb{N}$ なので、$\alpha$ が $\mathbb{N}$ の上界であることに
矛盾する。
\end{proof}

\begin{theorem}[有理数の稠密性]
$a<b$ ならば、ある $q\in\mathbb{Q}$ が存在して
\[
a<q<b
\]
となる。
\end{theorem}

\begin{proof}
アルキメデスの性質より、$n(b-a)>1$ となる自然数 $n$ を取る。整数 $m$ を
\[
m-1\le na<m
\]
となるように取る。このとき $a<m/n$ である。また
\[
\frac{m}{n}\le a+\frac{1}{n}<a+(b-a)=b
\]
である。よって $q=m/n$ とすればよい。
\end{proof}

\section*{3. 数列の極限}

\begin{definition}
数列 $\{a_n\}$ が実数 $a$ に収束するとは、任意の $\varepsilon>0$ に対して
ある $N\in\mathbb{N}$ が存在し、$n>N$ ならば
\[
|a_n-a|<\varepsilon
\]
となることをいう。このとき
\[
\lim_{n\to\infty}a_n=a,\qquad a_n\to a
\]
と書く。
\end{definition}

\begin{example}
\[
\frac{1}{2n+1}\to 0
\]
である。実際、$\varepsilon>0$ に対して $N>1/(2\varepsilon)$ となる
$N\in\mathbb{N}$ を取れば、$n>N$ のとき
\[
\left|\frac{1}{2n+1}-0\right|<\frac{1}{2n}<\varepsilon
\]
となる。
\end{example}

\begin{example}
\[
\frac{n^2-n}{n^2+1}\to 1
\]
である。実際、
\[
\left|\frac{n^2-n}{n^2+1}-1\right|
=\frac{n+1}{n^2+1}\le \frac{2}{n}
\]
なので、$N>2/\varepsilon$ とすればよい。
\end{example}

\section*{4. 極限の一意性}

\begin{theorem}
$a_n\to a$ かつ $a_n\to b$ ならば $a=b$ である。
\end{theorem}

\begin{proof}
任意の $\varepsilon>0$ を取る。$a_n\to a$ より、ある $N_1$ が存在して
$n>N_1$ ならば $|a_n-a|<\varepsilon/2$ である。同様に $a_n\to b$ より、
ある $N_2$ が存在して $n>N_2$ ならば $|a_n-b|<\varepsilon/2$ である。
$N=\max\{N_1,N_2\}$ とし、$n>N$ とすれば
\[
|a-b|\le |a-a_n|+|a_n-b|<\varepsilon
\]
となる。任意の $\varepsilon>0$ に対して成り立つので、$a=b$ である。
\end{proof}

\pointbox{
$\varepsilon$-$N$ 論法では、最後に示したい不等式から逆算して、
どの程度 $n$ を大きくすればよいかを決める。
}

\ifPDFTeX
\end{CJK}
\fi
\end{document}
